\begin{proposition}[P4.20]\hypertarget{tag-02WS}{}\label{night:P4.20}\label{stmt:P4.20}
\StatementRecord{night:P4.20}{proved}{02WS}{\SeriesRef{night:P4.16}, \SeriesRef{night:P4.6}, \SeriesRef{night:P4.8}}

Let $mS$ be an irreducible fibre of a proper holomorphic
degeneration of complex two-tori with smooth total threefold.
Suppose the normalization of $S$ is smooth and, outside
finitely many points, $S$ is smooth or has the ordinary-double
germ $xy=0$. Then $S$ has normal crossings everywhere.
In the global setting of 4.2, P4.6--P4.8 apply to this subclass.
\end{proposition}
\begin{proof}

\begin{enumerate}[label={$\langle1\rangle$\arabic*.},leftmargin=*]
\item We give the boundary and transport arguments explicitly.
Take a small ball $B$ around an exceptional point, with
boundary transverse to the smooth and double strata away
from that point. Let $g$ be a reduced equation. An analytic
$m$-th root of the unit in the actual defining function lets
us write that function as $g^m$.

\item Consider a small holomorphic deformation $G_s$ of $g$ whose
zero set near $\partial B$ retains its smoothly moving smooth
and ordinary-double strata, transverse to $\partial B$.
There is a uniform $\delta>0$ such that
$G_s|_{\partial B}$ is submersive over $0<|t|<\delta$.
At a boundary point on a double curve, choose coordinates
$a,b,c$ with $G_s=ab$, absorbing a unit in $a$.
Transversality of the double curve to $\partial B$ allows one
real $c$-coordinate to be solved from the boundary equation.
The complex coordinates $a,b$ remain free there, so $ab$ is
submersive when nonzero. At a smooth point use $G_s=a$.
Compactness gives uniform bounds; points outside these
neighbourhoods stay away from zero.

\item For each fixed small $s$ there are finitely many nonzero
critical values in this disc. After shrinking, the analytic
critical set in a neighbourhood of the compact closed ball
has finitely many relevant local components. The function is constant on
each analytic critical component. A component at a nonzero
small value cannot meet $\partial B$, so it is compact;
the maximum principle excludes a positive-dimensional
compact analytic subset of a ball. Thus the relevant critical
points are isolated and finite. Paths avoiding these values
join a small regular value to one fixed regular value $t_0$.
Proper Ehresmann transport on the closed ball, with the
verified boundary submersion, identifies the fibres along
such paths. For sufficiently small deformation the fibres
$G_s=t_0$ also form a smooth proper family to $s=0$, by
compactness. Hence every compact spherical class of nonzero
square in a local nearby fibre of $G_s$ transports into
$g=t_0$. Its inclusion in the actual torus fibre is impossible.

\item At a normalization preimage write the map germ as
$\phi:(\mathbb C^2,0)\to(\mathbb C^3,0)$. Away from the
exceptional point it is immersive with transverse multiple
images. If $d\phi$ drops rank, its maximal minors define a
nonempty finite scheme. Small generic polynomial
perturbations have transverse rank-one points near it.
We check both genericity and conservation of their positive
number. Varying linear coefficients makes the universal
derivative map submersive onto the $3$-by-$2$ matrix space.
The rank-one locus is smooth of codimension two; the rank-zero
locus has codimension six. Generic perturbations are transverse
to the first and avoid the second. The height-two maximal-minor
ideal has its Hilbert--Burch resolution. Its one-parameter
zero scheme is Cohen--Macaulay of dimension one; the finite
special fibre excludes a vertical component. Thus it is flat
and preserves the positive length. Further generic polynomial
coefficients keep the finitely many rank-one images off other
sheets: incidence of a zero-dimensional critical set with a
two-dimensional image in a three-dimensional target has
negative expected dimension.

\item Each transverse rank-one point is a Whitney crosscap, as we
now verify. Choose its first target coordinate as a source
coordinate $u$, and write $\phi=(u,F(u,v),H(u,v))$.
Transversality of the equations $F_v=H_v=0$ gives a linear
combination with nonzero second $v$-derivative. The
holomorphic Morse lemma with parameter $u$, followed by
subtracting the critical value, makes the second coordinate
$v^2$. Subtract the even part of the third coordinate to
write it as $vA(u,v^2)$. Transversality gives $A_u(0,0)\ne0$.
The target change $u\mapsto A(u,Y)$ and the matching source
change give exactly $(u,v)\mapsto(u,v^2,uv)$.
Its image is the crosscap of P4.16, with a compact spherical
class of square $-4$ in its nearby fibre.

\item The deformed images have holomorphic equations $G_s$ with
$G_0=g$. To see this, a generic projection of the original
finite normalization to two target coordinates is finite;
Weierstrass preparation preserves this after shrinking the
deformation. The total image is an analytic hypersurface
in $\mathbb C^3\times\Delta$. Its central equation has
multiplicity one, as seen at a smooth point of $S$; a unit
adjustment gives $G_0=g$. Outside a smaller ball the original
map is a transverse immersion, and this compact boundary
data persist under small deformation. We therefore have
the boundary property already proved. Transporting the
crosscap class contradicts its torus inclusion. The
normalization is consequently immersive everywhere.

\item Write its image germ as a unit times $g_1\cdots g_r$, where
each branch is smooth and $v_i=dg_i(0)\ne0$. Absorb the
unit in one factor. Independent gradients give $r\leq3$
and normal-crossing coordinates by the inverse function
theorem. Suppose instead $\sum_i c_iv_i=0$ is a nonzero
relation, and let $I$ be its nonzero support. Set
$a_i=\varepsilon/c_i$ for $i\in I$ and $a_j=\eta$ for
$j\notin I$, where $0<|\varepsilon|\ll|\eta|\ll1$.
Perturb constants so that $h_i(0)=a_i$. Keep all gradients
except one $i_0\in I$, and change that gradient by
\[
 \Delta v_{i_0}=-a_{i_0}\sum_{j\notin I}\frac{v_j}{a_j}.
\]
This is arbitrarily small and makes $\sum_i dh_i(0)/a_i=0$.
Thus $H=\prod_i h_i$ is critical at the origin and has
nonzero critical value $\prod_i a_i$. Its Hessian divided
by that value is
\[
 \sum_i\frac{\operatorname{Hess}(h_i)}{a_i}
 -\sum_i\frac{dh_i\otimes dh_i}{a_i^2}.
\]
An arbitrarily small generic quadratic change to $h_{i_0}$
makes this matrix invertible: its determinant is a nonzero
polynomial in the freely variable symmetric Hessian.
The branch functions remain smooth and retain their
ordinary-double/smooth boundary strata. We have constructed
a holomorphic Morse point off the divisor.

\item Its nearby fibre contains the real vanishing sphere in
$w_1^2+w_2^2+w_3^2=\epsilon$. Its normal bundle is
$iTS^2$, with the opposite orientation to $TS^2$ under the
complex surface orientation. Its self-intersection is
therefore $-e(S^2)=-2$. Boundary transport again puts this
class in the original torus fibre, a contradiction. Hence
the original branch gradients are independent everywhere,
which proves the normal-crossing assertion.

\end{enumerate}
\end{proof}